% Materiał z AF 25; fragment do rozdziału o splocie.
\subsection{Przestrzeń sprzężona do funkcji gładkich na zbiorze otwartym}
\begin{twierdzenie}[Dystrybucje o zwartym nośniku]
\label{thm:compact-support-E25}
Niech $\Omega\subseteq\R^d$ będzie niepustym zbiorem otwartym.
Obraz kanonicznego odwzorowania
$\mathcal E'(\Omega)\to\mathcal D'(\Omega)$
składa się dokładnie z dystrybucji o zwartym nośniku.
Dokładniej, dla każdego $F\in\mathcal E'(\Omega)$ istnieją
niepusty zwarty $K\Subset\Omega$, $m\in\N_0$ i $C>0$ takie, że
\[
 |F(f)|\le C\max_{|\alpha|\le m}\sup_{x\in K}
 |\partial^\alpha f(x)|\qquad(f\in\mathcal E(\Omega)).
\]
Odwrotnie, jeżeli $T\in\mathcal D'(\Omega)$ ma zwarty nośnik,
to istnieje dokładnie jeden funkcjonał ciągły
$F\in\mathcal E'(\Omega)$ równy $T$ na $\mathcal D(\Omega)$.
\end{twierdzenie}
\begin{proof}
Z kryterium ciągłości funkcjonału na przestrzeni
$\mathcal E(\Omega)$ (Twierdzenie~\ref{thm:dual-cinf17})
otrzymujemy oszacowanie przez skończenie wiele
seminorm $p_{K_j,m_j}$. Biorąc $K=\bigcup_jK_j$
i $m=\max_jm_j$, otrzymujemy oszacowanie z tezy
(dla funkcjonału zerowego można wybrać dowolny
niepusty zwarty $K$). Ograniczenie
$T=F|_{\mathcal D(\Omega)}$ jest dystrybucją, ponieważ
odwzorowanie $\mathcal D(\Omega)\to\mathcal E(\Omega)$
jest ciągłe na mocy twierdzenia~\ref{thm:zanurzenia25}. Gdy
$\varphi\in\mathcal D(\Omega\setminus K)$, jej
przedłużenie przez zero wraz ze wszystkimi pochodnymi
znika w otoczeniu $K$. Zatem $T(\varphi)=0$,
czyli $\supp T\subseteq K$.
Nośnik jest tu nośnikiem dystrybucji $T$,
a nie zbiorem punktów przestrzeni argumentów funkcjonału $F$.

Niech teraz $K=\supp T\Subset\Omega$. Jeśli $K$ jest
pusty, to $T=0$ z lokalności dystrybucji
(Lemat~\ref{lem:distribution-local19}) i wystarcza $F=0$.
W przeciwnym przypadku wybierzmy
$\chi\in\mathcal D(\Omega)$ równą $1$ w otoczeniu $K$
(Lemat~\ref{lem:cutoffs17}) i połóżmy $L=\supp\chi$.
Kładziemy $F(f)=\langle T,\chi f\rangle$ dla
$f\in\mathcal E(\Omega)$. Ponieważ $\supp(\chi f)\subseteq L$,
z lokalnego oszacowania dystrybucji
(Twierdzenie~\ref{thm:distribution-estimate19})
istnieją $C_L>0$ i $m\in\N_0$ takie, że
\[
 |F(f)|\le C_L p_{L,m}(\chi f)
 \le C_L 2^m p_{L,m}(\chi)p_{L,m}(f).
\]
Ostatnia nierówność wynika z reguły Leibniza i tożsamości
$\sum_{\gamma\le\alpha}\binom{\alpha}{\gamma}=2^{|\alpha|}$.
Stąd $F$ jest liniowym funkcjonałem ciągłym.
Wartość nie zależy od wyboru $\chi$: dla dwóch odcięć
$\chi_1,\chi_2$ funkcja $(\chi_1-\chi_2)f$ jest testowa
i znika w otoczeniu $K$. Jej nośnik jest więc zwartym
podzbiorem $\Omega\setminus K$. Z definicji nośnika
i lokalności dystrybucji wynika, że $T$ ją anihiluje. Jeżeli $f\in\mathcal D(\Omega)$, to
$T(f)=T(\chi f)$, ponieważ $(1-\chi)f$ znika w
otoczeniu $K$. Jednoznaczność przedłużenia wynika
z gęstości $\mathcal D(\Omega)$ w $\mathcal E(\Omega)$.
Otrzymana identyfikacja jest zgodna z przedłużeniem
z Lematu~\ref{lem:compact-action21} oraz Uwagą~
\ref{rem:Eprime21}; nie nadajemy obrazowi automatycznie
topologii podprzestrzeni $\mathcal D'$.
\end{proof}
